% ============================================================================
% Fichier  : partie3/P03_C09_normes_internationales.tex
% Titre    : Les Normes Internationales de la Gouvernance et des Services
% Auteur   : MINKA MI NGUIDJOÏ Thierry Emmanuel
% Version  : 1.0 -- octobre 2026
% ============================================================================
% Sources : normes du dossier Normes (ISO/IEC 27000:2018, 27001:2013,
% 27002:2013, 27006 ; résumé COSO ERM 2004 en français) ; éditions en
% vigueur citées d'après les catalogues ISO. Principes de l'ISO/IEC
% 38500:2024 repris d'ISO 37000 : liste à vérifier dans le texte de la
% norme.
% ============================================================================

\chapter{Les Normes Internationales de la Gouvernance et des Services}
\label{chap:normes}
\chaptermark{Normes internationales}

\epigraph{%
    « Une norme ne dit pas ce qu'il faut penser.\\
    Elle dit ce qu'il faudra pouvoir démontrer. »%
}{Adage des auditeurs de systèmes de management}

\niveauChapitre{Fondamental}{%
    Maîtrise du Chapitre~\ref{chap:referentiels} (natures des référentiels
    et carte d'articulation).%
}

\begin{boxobjectifs}
\begin{itemize}
    \item Présenter les principes et le modèle « évaluer, diriger,
          surveiller » de l'ISO/IEC~38500.
    \item Décrire les exigences d'un système de management des services
          (ISO/IEC~20000-1).
    \item Articuler la gouvernance des SI avec l'ISO/IEC~27001, l'ISO~31000
          et COSO ERM.
    \item Distinguer certification, conformité et alignement volontaire.
\end{itemize}
\end{boxobjectifs}

\bigskip

% ============================================================================
\section*{Question directrice : faut-il un certificat~?}
% ============================================================================

KRIBILOG SA vient de recevoir l'appel d'offres d'un grand armateur pour
la gestion de ses conteneurs au port de Kribi. Une clause l'arrête~: « Le
prestataire devra justifier, dans un délai de dix-huit mois, d'une
certification ISO/IEC~20000-1 sur le périmètre du service. » Au même
moment, le conseil de la BANQUE SANAGA demande s'il ne faudrait pas
« se faire certifier ISO~27001 pour rassurer la COBAC ».

La Partie~\ref{part:normatif}, qui s'ouvre ici, place la gouvernance des SI
dans son cadre normatif, puis juridique (Chapitre~\ref{chap:droit}), puis
dans celui de l'audit (Chapitre~\ref{chap:audit}). Ce chapitre présente les
normes internationales qui comptent pour la gouvernance des SI~: leur
architecture commune (section~\ref{sec:c09-structure}), la norme de la
gouvernance (section~\ref{sec:c09-38500}), celles des services et de la
sécurité (sections~\ref{sec:c09-20000} et~\ref{sec:c09-27001}), celles du
risque (section~\ref{sec:c09-risque}), et enfin la question du certificat
(section~\ref{sec:c09-certification}).

\separateur

% ============================================================================
\section{L'Architecture des Systèmes de Management}
\label{sec:c09-structure}
% ============================================================================

\subsection{Qu'est-ce qu'un système de management~?}

\begin{boxdef}
\termecle{Système de management}\index{système de management}
(\textit{management system})~: ensemble d'éléments corrélés ou en
interaction d'une organisation (politiques, objectifs, rôles, processus)
qui sert à établir des politiques et des objectifs, et à atteindre ces
objectifs. Un système de management de la qualité, de la sécurité de
l'information ou des services en sont des exemples.
\end{boxdef}

\subsection{Une structure harmonisée}

Les normes ISO de systèmes de management certifiables partagent une même
structure, dite \termecle{structure harmonisée}\index{structure
harmonisée} (anciennement « structure de haut niveau »). Elle facilite la
mise en place de systèmes intégrés~: une organisation certifiée à la fois
ISO/IEC~27001 et ISO/IEC~20000-1 n'a qu'un seul cycle de revue de
direction, d'audit interne et d'amélioration.

\begin{tableaucours}{cL{3.6cm}L{8.0cm}}{%
    \tableauHeader{Art.} &
    \tableauHeader{Titre} &
    \tableauHeader{Contenu}
}
4 & Contexte de l'organisation & Enjeux internes et externes, parties
intéressées, périmètre du système \\
5 & Leadership & Engagement de la direction, politique, rôles et
responsabilités \\
6 & Planification & Risques et opportunités, objectifs et plans pour les
atteindre \\
7 & Support & Ressources, compétences, sensibilisation, communication,
informations documentées \\
8 & Fonctionnement & Planification et maîtrise opérationnelles~: c'est là
que chaque norme place ses exigences propres \\
9 & Évaluation des performances & Surveillance, mesure, audit interne,
revue de direction \\
10 & Amélioration & Non-conformités, actions correctives, amélioration
continue \\
\end{tableaucours}
\vspace{-0.8em}
\begin{center}{\small\itshape Tableau~9.1. La structure harmonisée des
normes ISO de systèmes de management (articles~4 à~10)}\end{center}

Cette structure suit le cycle \emph{planifier, réaliser, vérifier, agir}~:
les articles~4 à~7 planifient, l'article~8 réalise, l'article~9 vérifie,
l'article~10 améliore. On y retrouve, à une autre échelle, la logique de
COBIT~: le domaine EDM correspond aux articles~5 et~9 (direction et
surveillance), les domaines de management aux articles~6 à~8 et~10.

% ============================================================================
\section{ISO/IEC 38500 : la Norme de la Gouvernance}
\label{sec:c09-38500}
% ============================================================================

\subsection{Une norme pour l'organe de gouvernance}

L'ISO/IEC~38500 s'adresse à l'organe de gouvernance~: conseil
d'administration, conseil de surveillance, direction d'un établissement
public. Elle ne décrit ni processus ni contrôles, mais des principes et un
modèle. Elle n'est pas certifiable. Sa troisième édition a été publiée en
février~2024~\cite{iso38500_2024}.

\subsection{Le modèle évaluer, diriger, surveiller}

\noindent\begin{minipage}{\linewidth}
\begin{center}
\begin{tikzpicture}[
    tache/.style = {
        draw=CouleurPrimaire, fill=CouleurPrimaire, rounded corners=4pt,
        text width=2.4cm, minimum height=0.9cm, align=center,
        font=\footnotesize\bfseries\color{white}
    },
    bloc/.style = {
        draw=CouleurBordDef, fill=CouleurDef, rounded corners=4pt,
        text width=3.0cm, minimum height=0.9cm, align=center,
        font=\footnotesize\bfseries\color{CouleurPrimaire}
    },
    fleche/.style = {-{Stealth[length=5pt]}, thick, color=CouleurSecondaire}
]
\node[tache] (ev) at (0,2.2)   {Évaluer};
\node[tache] (di) at (3.6,2.2) {Diriger};
\node[tache] (su) at (7.2,2.2) {Surveiller};
\node[bloc] (pr) at (1.8,0)    {Propositions,\\ plans};
\node[bloc] (op) at (5.4,0)    {Activités,\\ performance};
\draw[fleche] (ev) -- (di);
\draw[fleche] (di) -- node[left, font=\scriptsize\itshape, align=right]{politiques,\\
    plans} (op);
\draw[fleche] (op) -- node[right, font=\scriptsize\itshape]{mesures} (su);
\draw[fleche] (pr) -- (ev);
\draw[fleche] (su.north) to[out=120, in=60] (ev.north);
\node[font=\scriptsize\bfseries\color{CouleurAccent}] at (3.6,4.0)
    {Organe de gouvernance};
\node[font=\scriptsize\bfseries\color{CouleurAccent}] at (3.6,-0.9)
    {Management};
\end{tikzpicture}
\end{center}
\vspace{-0.4em}
\begin{center}{\small\itshape Figure~9.1. Le modèle évaluer, diriger,
surveiller (représentation simplifiée)}\end{center}
\end{minipage}

\medskip

L'organe de gouvernance \textbf{évalue} l'usage actuel et futur de
l'informatique, ainsi que les propositions que lui soumet le management~;
il \textbf{dirige} la préparation et la mise en œuvre des plans et des
politiques~; il \textbf{surveille} la performance et la conformité. Ce sont
les trois verbes du domaine EDM de COBIT, qui s'inspire directement de la
norme.

\subsection{Des principes de 2015 à ceux de 2024}

L'édition de 2015~\cite{iso38500} reposait sur six principes~:
\emph{responsabilité}, \emph{stratégie}, \emph{acquisition},
\emph{performance}, \emph{conformité} et \emph{comportement humain}.
L'édition de 2024 les remplace par les onze principes de la norme
ISO~37000 sur la gouvernance des organisations, et précise comment ils
s'appliquent à la gouvernance des SI. Elle fait de la gouvernance des SI
une partie intégrante de la gouvernance de l'organisation, et l'oriente
vers trois résultats~: une performance effective, une gestion
responsable et un comportement éthique.

\begin{tableaucours}{L{3.2cm}L{9.0cm}}{%
    \tableauHeader{Catégorie} &
    \tableauHeader{Principes (d'après ISO~37000)}
}
Principe premier & Raison d'être (\textit{purpose}) \\
Principes fondateurs & Création de valeur~; stratégie~; supervision~;
redevabilité \\
Principes facilitateurs & Engagement des parties prenantes~; leadership~;
données et décisions~; gouvernance des risques~; responsabilité
sociétale~; viabilité et performance dans la durée \\
\end{tableaucours}
\vspace{-0.8em}
\begin{center}{\small\itshape Tableau~9.2. Les onze principes de
gouvernance repris par l'ISO/IEC~38500:2024 (notre traduction)}\end{center}
% A VERIFIER dans le texte de l'ISO/IEC 38500:2024 : intitulés exacts et
% ordre des onze principes (source : présentations publiques d'ISO 37000).

\begin{boxscenario}{Opération MEFOU. Une charte inspirée de l'ISO/IEC 38500}
Le nouveau président du conseil d'administration de l'ANPP demande que la
charte du comité stratégique des SI (Chapitre~\ref{chap:gouv-entreprise})
s'appuie sur l'ISO/IEC~38500. La division informatique propose trois
articles~:
\begin{enumerate}
    \item \emph{Raison d'être}~: tout projet informatique doit être rattaché
          à la mission de l'agence (payer l'État à bonne date) et à un
          programme budgétaire.
    \item \emph{Supervision}~: le comité \emph{évalue} chaque année le
          portefeuille, \emph{dirige} par la décision d'investissement et
          \emph{surveille} chaque trimestre les projets en cours.
    \item \emph{Données et décisions}~: le conseil reçoit, avant toute
          décision d'investissement, un dossier normalisé (coût total,
          bénéfices, risques).
\end{enumerate}
La charte ne cite pas la norme pour la forme~: elle en traduit les
principes en règles vérifiables.
\end{boxscenario}

% ============================================================================
\section{ISO/IEC 20000-1 : le Management des Services}
\label{sec:c09-20000}
% ============================================================================

\subsection{Les exigences}

L'ISO/IEC~20000-1, dont l'édition en vigueur date de 2018, formule les
exigences d'un \termecle{système de management des services}\index{système
de management des services} (\sigle{SMS}). Elle suit la structure
harmonisée, et place à l'article~8 ses exigences propres au
fonctionnement des services.

\begin{tableaucours}{cL{3.6cm}L{8.0cm}}{%
    \tableauHeader{Art.} &
    \tableauHeader{Thème} &
    \tableauHeader{Exigences principales}
}
8.2 & Portefeuille de services & Planifier les services, gérer le catalogue,
les actifs et la configuration \\
8.3 & Relations et accords & Relations avec les clients, gestion des
niveaux de service, gestion des fournisseurs \\
8.4 & Offre et demande & Budgétisation des services, demande et capacité \\
8.5 & Conception, construction et transition & Gestion des changements,
conception et transition des services, mises en production et
déploiement \\
8.6 & Résolution et réalisation & Gestion des incidents, des demandes de
service et des problèmes \\
8.7 & Assurance des services & Disponibilité, continuité, sécurité de
l'information \\
\end{tableaucours}
\vspace{-0.8em}
\begin{center}{\small\itshape Tableau~9.3. Les exigences de l'article~8 de
l'ISO/IEC~20000-1:2018 (synthèse, notre traduction)}\end{center}

\subsection{ITIL, COBIT et ISO/IEC 20000-1}

Les trois référentiels se recouvrent largement sur la gestion des services,
et c'est précisément pourquoi il faut savoir ce que chacun apporte.
\textbf{ITIL} décrit \emph{comment} mener chaque pratique~;
\textbf{ISO/IEC~20000-1} dit \emph{ce qui doit être démontré} pour obtenir
un certificat~; \textbf{COBIT} relie ces pratiques aux objectifs de
l'organisation et à sa gouvernance. Une organisation qui vise la
certification ISO/IEC~20000-1 s'appuie le plus souvent sur ITIL pour
construire ses pratiques. Le Chapitre~\ref{chap:processus} établira la
correspondance détaillée entre les trois.

\begin{boxmetier}{Responsable des services IT}
\textbf{Préparer KRIBILOG à l'ISO/IEC 20000-1}
\begin{enumerate}
    \item Définir le périmètre~: le service de suivi des conteneurs, depuis
          Douala et Kribi, et non toute l'informatique de l'entreprise.
    \item Établir que KRIBILOG reste \emph{responsable} du service même s'il
          est exploité par NETPORT SERVICES~: la norme exige la maîtrise
          des processus confiés à des tiers.
    \item Mettre d'abord en place les pratiques de l'article~8.6
          (incidents, demandes, problèmes) et de l'article~8.5
          (changements), qui sont aujourd'hui les plus défaillantes.
    \item Conduire un audit interne et une revue de direction avant de
          solliciter l'organisme de certification.
\end{enumerate}
Dix-huit mois sont un délai serré, mais réaliste si le périmètre reste
limité au service demandé par l'armateur.
\end{boxmetier}

% ============================================================================
\section{ISO/IEC 27001 et 27002 : la Sécurité de l'Information}
\label{sec:c09-27001}
% ============================================================================

\subsection{Le système de management de la sécurité}

L'ISO/IEC~27001 formule les exigences d'un \termecle{système de management
de la sécurité de l'information}\index{système de management de la
sécurité de l'information} (\sigle{SMSI}). Le vocabulaire de la famille
est fixé par l'ISO/IEC~27000~\cite{iso27000_2018}. La démarche est fondée
sur le risque~: l'organisation apprécie ses risques de sécurité (selon une
méthode que précise l'ISO/IEC~27005), choisit les mesures qui les
traitent, et justifie ses choix.

\begin{boxdef}
\termecle{Déclaration d'applicabilité}\index{déclaration d'applicabilité}
(\textit{statement of applicability}, \sigle{SoA})~: document qui liste les
mesures de sécurité de l'annexe~A de l'ISO/IEC~27001, indique pour
chacune si elle est retenue ou exclue, et justifie ce choix. C'est la
pièce centrale d'un audit de certification ISO/IEC~27001.
\end{boxdef}

\subsection{Deux éditions à ne pas confondre}

\begin{tableaucours}{L{3.4cm}L{4.4cm}L{4.4cm}}{%
    \tableauHeader{Élément} &
    \tableauHeader{Édition 2013} &
    \tableauHeader{Édition 2022}
}
Mesures de l'annexe~A & 114~mesures en 14~domaines & 93~mesures en
4~thèmes \\
Thèmes & · & Organisationnelles (37), humaines (8), physiques (14),
technologiques (34) \\
Nouvelles mesures & · & Notamment renseignement sur les menaces,
sécurité de l'informatique en nuage, préparation des TIC à la continuité \\
Statut & Remplacée & En vigueur \\
\end{tableaucours}
\vspace{-0.8em}
\begin{center}{\small\itshape Tableau~9.4. ISO/IEC~27001 et 27002~: éditions
2013 et 2022}\end{center}

\begin{boxalert}
\textbf{Violation typique.} Une organisation rédige en 2026 sa politique de
sécurité et sa déclaration d'applicabilité à partir des éditions 2013 des
normes, trouvées dans un dossier partagé. L'auditeur de certification
relève aussitôt l'écart~: la référence n'est plus celle en vigueur, et les
mesures relatives à l'informatique en nuage manquent.

\smallskip
\textbf{Règle~:} toujours vérifier l'édition en vigueur d'une norme dans le
catalogue de l'ISO avant de s'en servir de référence, et indiquer l'édition
dans chaque document.
\end{boxalert}

\subsection{Le lien avec la gouvernance}

La sécurité de l'information n'est pas la gouvernance des SI
(Chapitre~\ref{chap:gouv-entreprise}), mais elle en relève par trois
voies. Le conseil fixe l'appétence au risque de sécurité (EDM03). Le
management conduit le SMSI (APO13, dont la finalité recouvre celle de
l'ISO/IEC~27001). L'exploitation met en œuvre les mesures (DSS05). Un
SMSI certifié est un excellent moyen d'atteindre APO13~; il ne dit rien
de la valeur des investissements ni de la qualité des services.

\begin{boxmetier}{RSSI}
\textbf{Répondre au conseil de la BANQUE SANAGA}
\begin{enumerate}
    \item Une certification ISO/IEC~27001 ne « rassure » le régulateur que
          si son périmètre couvre les systèmes qui l'intéressent~: les
          systèmes de paiement et le système central.
    \item Elle ne remplace ni les audits prévus par la loi n°~2010/012
          (Chapitre~\ref{chap:referentiels}), ni les exigences de la COBAC
          sur le contrôle interne.
    \item Elle n'aurait pas évité la panne de mars~2026, qui ne relevait
          pas de la sécurité mais de la gestion des changements.
    \item Recommandation~: viser la certification à trois ans sur le
          périmètre des paiements, comme le prévoit la carte des
          référentiels (Chapitre~\ref{chap:referentiels}), et traiter
          d'abord la feuille de route COBIT (Chapitre~\ref{chap:cobit-capacite}).
\end{enumerate}
\end{boxmetier}

% ============================================================================
\section{Le Risque : ISO 31000 et COSO}
\label{sec:c09-risque}
% ============================================================================

\subsection{ISO 31000 : des lignes directrices pour tout risque}

L'ISO~31000 (édition 2018) donne des lignes directrices sur le management
du risque, applicables à tout type de risque et à toute organisation. Elle
n'est pas certifiable. Elle s'articule en trois parties~: des
\emph{principes} (un management du risque intégré, structuré, adapté,
inclusif, dynamique, fondé sur la meilleure information disponible,
attentif aux facteurs humains et culturels, et en amélioration continue),
un \emph{cadre organisationnel}, et un \emph{processus}~: établir le
contexte et les critères, apprécier les risques (les identifier, les
analyser, les évaluer), les traiter, puis surveiller, communiquer et
rendre compte.

\subsection{COSO : contrôle interne et risques de l'entreprise}

Les deux cadres du COSO (Chapitre~\ref{chap:referentiels}) sont la
référence des auditeurs et des régulateurs financiers.

\begin{tableaucours}{L{3.6cm}L{8.6cm}}{%
    \tableauHeader{Cadre} &
    \tableauHeader{Composantes}
}
Contrôle interne (2013) & Environnement de contrôle~; évaluation des
risques~; activités de contrôle~; information et communication~;
activités de pilotage. Cinq composantes et 17~principes. \\
Management des risques (2004) & Huit éléments~: environnement interne~;
fixation des objectifs~; identification des événements~; évaluation des
risques~; traitement des risques~; activités de contrôle~; information et
communication~; pilotage~\cite{coso2004ermfr}. \\
Management des risques (2017) & Cinq composantes~: gouvernance et
culture~; stratégie et définition des objectifs~; performance~; revue et
révision~; information, communication et reporting. Vingt
principes~\cite{coso2017erm}. \\
\end{tableaucours}
\vspace{-0.8em}
\begin{center}{\small\itshape Tableau~9.5. Les cadres du COSO}\end{center}

Le résumé en français disponible dans le dossier \texttt{Normes} est celui
de l'édition de 2004. L'édition de 2017 en déplace le centre de gravité~:
elle ne traite plus le risque comme un sujet de contrôle, mais comme une
composante de la stratégie et de la performance. C'est la lecture que
retient aussi COBIT avec la finalité d'optimisation des risques
(Chapitre~\ref{chap:alignement}).

\subsection{Articuler les risques}

La règle d'articulation est la suivante~: \textbf{le risque informatique
est un risque de l'entreprise}. Le registre des risques liés aux SI
(APO12) doit alimenter le registre des risques de l'organisation, selon
la même échelle et les mêmes critères, et l'appétence au risque fixée par
le conseil (EDM03) s'applique aux deux. Une organisation qui tient deux
registres séparés, avec deux échelles, ne peut pas comparer un risque
informatique à un risque de crédit~: elle ne peut donc pas arbitrer.

% ============================================================================
\section{Certification, Conformité, Alignement}
\label{sec:c09-certification}
% ============================================================================

\subsection{Trois démarches}

\begin{boxdef}
\termecle{Certification}\index{certification}~: attestation, par un
organisme tiers accrédité, qu'un système de management est conforme aux
exigences d'une norme, sur un périmètre donné et pour une durée donnée.

\medskip
\termecle{Conformité}\index{conformité}~: respect d'exigences obligatoires
(lois, règlements, contrats). Elle ne se choisit pas~; elle se démontre.

\medskip
\termecle{Alignement volontaire}\index{alignement volontaire}~: choix
d'une organisation de s'inspirer d'un référentiel ou d'une norme sans
viser de certificat. C'est le cas de tout usage de COBIT, d'ITIL ou de
l'ISO~31000.
\end{boxdef}

\subsection{Le cycle de certification}

Une certification de système de management suit un cycle de trois ans~:
un audit initial en deux étapes (revue de la documentation et de l'état
de préparation, puis audit de la mise en œuvre), la délivrance du
certificat, des audits de surveillance au moins annuels, puis un audit de
renouvellement. Les organismes de certification doivent eux-mêmes
satisfaire aux exigences de l'ISO/IEC~17021-1, complétées pour la sécurité
de l'information par l'ISO/IEC~27006, et être accrédités par un organisme
d'accréditation. Un certificat délivré par un organisme non accrédité n'a
qu'une valeur commerciale.

\begin{boxprocedure}
\textbf{Préparer un audit de certification.} Avant de solliciter
l'organisme de certification, l'organisation doit pouvoir présenter~:
\begin{itemize}
    \item le périmètre du système, écrit et justifié~;
    \item la politique et les objectifs, approuvés par la direction~;
    \item l'appréciation des risques et, pour l'ISO/IEC~27001, la
          déclaration d'applicabilité~;
    \item les procédures exigées et les enregistrements de leur
          fonctionnement sur plusieurs mois~;
    \item au moins un cycle complet d'audit interne et de revue de
          direction~;
    \item le suivi des non-conformités relevées et des actions correctives.
\end{itemize}
\end{boxprocedure}

\begin{boxscenario}{Opération KRIBI. Le calendrier de la certification}
\textbf{Contrainte.} Certification ISO/IEC~20000-1 exigée dans dix-huit
mois sur le service de suivi des conteneurs.

\medskip
\textbf{Plan retenu par M.~MBARGA.} Mois 1 à 2~: périmètre, politique et
objectifs, analyse d'écart. Mois 3 à 9~: mise en place des pratiques
(incidents, demandes, problèmes, changements, niveaux de service,
fournisseurs), avec un avenant au contrat de NETPORT SERVICES. Mois 10 à
13~: fonctionnement et enregistrements. Mois 14~: audit interne. Mois 15~:
revue de direction et actions correctives. Mois 16 à 17~: audit de
certification en deux étapes. Mois 18~: certificat.

\medskip
\textbf{Analyse} (confiance moyenne). Le calendrier tient si l'avenant avec
l'infogérant est signé dans les trois premiers mois~: sans lui, KRIBILOG
ne pourra démontrer la maîtrise d'un service exploité par un tiers.
\end{boxscenario}

\begin{boxmetier}{Auditeur SI}
\textbf{Lire un certificat en cinq questions}
\begin{enumerate}
    \item Quelle norme, et quelle édition~?
    \item Quel périmètre exact (sites, services, systèmes)~?
    \item Quel organisme l'a délivré, et est-il accrédité~?
    \item Quelle date de validité, et le dernier audit de surveillance
          a-t-il eu lieu~?
    \item Le périmètre couvre-t-il le service qui vous intéresse~?
\end{enumerate}
\end{boxmetier}

% ============================================================================
\separateur
% ============================================================================

\begin{boxresume}
\textbf{Ce qu'il faut retenir de ce chapitre~:}
\begin{itemize}
    \item \textbf{Une structure harmonisée}~: les normes de systèmes de
          management partagent les articles~4 à~10 (contexte, leadership,
          planification, support, fonctionnement, évaluation,
          amélioration), ce qui permet des systèmes intégrés.
    \item \textbf{ISO/IEC~38500} s'adresse à l'organe de gouvernance~:
          évaluer, diriger, surveiller. Depuis 2024, elle reprend les onze
          principes de l'ISO~37000. Elle n'est pas certifiable.
    \item \textbf{ISO/IEC~20000-1} exige un système de management des
          services~; ITIL dit comment le construire, COBIT le relie à la
          gouvernance.
    \item \textbf{ISO/IEC~27001 et 27002}~: un SMSI fondé sur le risque, une
          déclaration d'applicabilité, 93~mesures depuis 2022. Vérifier
          toujours l'édition en vigueur.
    \item \textbf{ISO~31000 et COSO}~: le risque informatique est un risque
          de l'entreprise, à inscrire dans un registre et une échelle
          communs.
    \item \textbf{Certification, conformité, alignement}~: la certification
          atteste un périmètre, la conformité répond à une obligation,
          l'alignement est un choix.
\end{itemize}
\end{boxresume}

\bigskip

\begin{boxexo}[title={\ding{45}\quad Exercice~9.1. Quel article~? \niveaubase}]
Indiquez l'article de la structure harmonisée (4 à 10) auquel se rattache
chacun des éléments suivants.
\begin{enumerate}[(a)]
    \item La revue annuelle du système par la direction générale.
    \item La liste des parties intéressées et de leurs attentes.
    \item Le plan de formation des techniciens du centre de services.
    \item Le traitement d'une non-conformité relevée en audit interne.
    \item La procédure de gestion des incidents.
\end{enumerate}
\end{boxexo}

\begin{boxexo}[title={\ding{45}\quad Exercice~9.2. Deux registres, une échelle \niveaumoyen}]
La BANQUE SANAGA tient un registre des risques de crédit et de marché
(échelle de 1 à 4) et, à la DSI, un registre des risques informatiques
(échelle de 1 à 5, définie différemment).
\begin{enumerate}[(a)]
    \item Expliquez pourquoi cette situation empêche le conseil
          d'arbitrer.
    \item Proposez une échelle commune, avec des critères d'impact
          exprimés en FCFA et en durée d'indisponibilité.
    \item Reclassez sur cette échelle le risque « échec d'une migration du
          système central ».
\end{enumerate}
\end{boxexo}

\begin{boxexo}[title={\ding{45}\quad Exercice~9.3. Une stratégie normative pour NOUN PAY \niveauexpert}]
Les investisseurs de NOUN PAY exigent « des certifications
internationales ». Les moyens de l'entreprise ne permettent d'en viser
qu'une dans les deux ans.
\begin{enumerate}[(a)]
    \item Comparez l'intérêt d'une certification ISO/IEC~27001 et d'une
          certification ISO/IEC~20000-1 pour NOUN PAY.
    \item Proposez un périmètre réaliste pour la certification retenue.
    \item Indiquez les référentiels que l'entreprise suivra par simple
          alignement volontaire, et pourquoi.
    \item Rédigez la réponse de Mme~KAMGA aux investisseurs, en distinguant
          certification, conformité légale (loi n°~2024/017) et
          alignement.
\end{enumerate}
\end{boxexo}

\bigskip

\begin{boxinfo}
\textbf{Transition.} Les normes sont des engagements volontaires~; la loi,
elle, s'impose. Le Chapitre~\ref{chap:droit} présente le cadre juridique
et réglementaire camerounais et africain de la gouvernance des SI~:
cybersécurité, communications électroniques, données personnelles,
réglementation bancaire et droit des sociétés.
\end{boxinfo}
